\clearpage
\section*{Appendix AC. Local retention forces quadratic-stress concentration}
\addcontentsline{toc}{section}{Appendix AC. Local retention and stress concentration}
\noindent\textbf{Scope.} This appendix derives a source-independent necessary quadratic-stress gradient from pressure-driven inward energy transport and local retention. No separation or disjoint-support assumption is made on the nonlinear exterior correction. The statement is conditional on the nonpressure and carrier-only pressure fluxes being subcritical and does not construct an unforced singularity.

\paragraph{Theorem AC.1 (Pressure-to-stress concentration).}
Let $0<h<1/6$, $\tau=T-t$, $p=\tfrac12-3h>0$ and
$I_\tau=[T-\tau,T-\tau/2]$. Assume the smooth periodic divergence-free
unforced velocity $u=U+w$ and moving regions $\Omega(t)$ satisfy
all core and local-retention hypotheses of Appendix T, with
\[
 \|U(t)\|_2^2=M_0\tau^p(1+o(1)),\qquad
 \|\nabla U(t)\|_2^2=R_0M_0\tau^{p-1}(1+o(1)),
 \quad M_0>0,\quad R_0>p/2.
\]
Assume, in addition, the integrated sum of nonpressure boundary
flux and \emph{carrier-only} pressure flux is $o(\tau^p)$:
\begin{equation}\label{eq:AC-other}
\int_{I_\tau}\bigl[\Phi_{\rm other}(t)+
\int_{\partial\Omega(t)}\pi[U](u\cdot n)\,dS\bigr]dt
=o(\tau^p).
\end{equation}
The premise can follow from the additional interface bounds of
Appendix U and the periodic carrier-pressure collar limit of Appendix AA,
but it is not implied by local $L^2$ retention alone.

Define $\delta_\tau=\tau^{-p/2}\sup_{t\in I_\tau}
\|w(t)\|_{L^2(\Omega(t))}$, so $\delta_\tau\to0$,
and let
\[
 \pi_c=\pi[U+w]-\pi[U]=(-\Delta)^{-1}_0\partial_i\partial_j S_{ij},
 \qquad S=U\otimes w+w\otimes U+w\otimes w.
\]
Here $S$ is the \emph{velocity quadratic stress}, not the positive
seven-ray H8 engineering stress. For some $c_0>0$ and all sufficiently
small $\tau$,
\begin{equation}\label{eq:AC-L2}
\boxed{
\sup_{t\in I_\tau}\|\nabla\pi_c(t)\|_{L^2(\Omega(t))}
\ \ge\ c_0\,\delta_\tau^{-1}\tau^{p/2-1}.}
\end{equation}
Moreover, with the globally Fourier-defined $L^2$ tensor norm,
\begin{equation}\label{eq:AC-stress}
\boxed{
\sup_{t\in I_\tau}\|\nabla S(t)\|_{L^2(\mathbb T^3)}
\ \ge\ c_0\,\delta_\tau^{-1}\tau^{p/2-1},
\qquad
\tau^{1-p/2}\sup_{t\in I_\tau}\|\nabla S(t)\|_2\longrightarrow\infty.
}
\end{equation}
If $\delta_\tau=0$ on any sufficiently late dyad, the assumed
inward pressure transfer is impossible.

\begin{proof}
Appendix T yields the dyadic inward energy estimate
\[
-\int_{I_\tau}\Phi(t)\,dt\ge(K+o(1))\tau^p,\qquad
K=\frac{1-2^{-p}}p(R_0-p/2)M_0>0.
\]
Subtracting \eqref{eq:AC-other} gives
$-\int_{I_\tau}\int_{\partial\Omega}\pi_c u\cdot n\,dSdt
\ge (K+o(1))\tau^p$. Since $U$ is solenoidal and vanishes on a
collar of the boundary, $\int_{\Omega(t)}U\cdot\nabla\pi_c=0$.
The exact divergence theorem therefore gives
\[
\int_{\partial\Omega(t)}\pi_c u\cdot n\,dS
=\int_{\Omega(t)}w\cdot\nabla\pi_c\,dx.
\]
By Cauchy--Schwarz, the time interval length $|I_\tau|=\tau/2$,
and the definition of $\delta_\tau$,
\[
 (K+o(1))\tau^p
 \le\frac\tau2\,\delta_\tau\tau^{p/2}
 \sup_{t\in I_\tau}\|\nabla\pi_c(t)\|_{L^2(\Omega(t))}.
\]
For small enough $\tau$, take $c_0=K$, yielding \eqref{eq:AC-L2}.

For $k\ne0$, the normalized periodic Poisson multiplier obeys
\[
 \widehat{\partial_\ell\pi_c}(k)
 =-ik_\ell\frac{k_i k_j}{|k|^2}\widehat S_{ij}(k).
\]
The matrix $(k_i k_j/|k|^2)_{ij}$ has Frobenius norm one.
Parseval and Cauchy--Schwarz give the exact universal inequality
\[
 \|\nabla\pi_c\|_{L^2(\mathbb T^3)}
 \le\|\nabla S\|_{L^2(\mathbb T^3)}.
\]
Its left side dominates the local $L^2(\Omega(t))$ norm.
Combining with \eqref{eq:AC-L2} proves \eqref{eq:AC-stress}.
\end{proof}

\paragraph{Corollary AC.2 (No-go for natural-scale quadratic stress).}
Under the hypotheses of Theorem AC.1, if for all late dyads
\[
\sup_{t\in I_\tau}
\|\nabla(U\otimes w+w\otimes U+w\otimes w)(t)\|_{L^2(\mathbb T^3)}
\le C\tau^{p/2-1},
\]
then the locally retained unforced solution cannot exist. No
spatial separation, fourfold symmetry, stress positivity, or
restriction on globally non-negligible exterior energy is assumed.
The critical growth exponent is $p/2-1=-3/4-3h/2$, equal to $-9/10$
when $h=1/10$.

\paragraph{Corollary AC.3 (Pointwise pressure concentration).}
For $\Omega(t)=x_0+\operatorname{diag}(\tau^{1/2},
\tau^{1/2},\tau^{1/2-h})\Omega_0$, the moving volume equals
$|\Omega_0|\tau^{3/2-h}$. Therefore
\begin{equation}\label{eq:AC-Linfty}
\boxed{\sup_{t\in I_\tau}\|\nabla\pi_c(t)\|_{L^\infty(\Omega(t))}
\ge\frac{c_0}{|\Omega_0|^{1/2}}\,
\delta_\tau^{-1}\tau^{-3/2-h}.}
\end{equation}
For the regular near-correction scaling $\delta_\tau=O(\tau^h)$,
this recovers or strengthens the $\tau^{-3/2-2h}$
gradient threshold of Appendix AB. The new inequality needs no
a priori $L^\infty$ bound on $w$; instead it explicitly retains the
local $L^2$ deficit.

\paragraph{Final claim boundary.}
Theorem AC.1 places a quantitative necessary condition on a
\emph{hypothetical} pressure-driven exact solution; its hypotheses
do not follow from the unforced equation without the remaining
source-specific checks. A correction with sufficiently concentrated
quadratic stress is not excluded by this result and may still fail
the positivity, source-matching, nonlinear zero-residual, and
terminal-retention gates. The original Case~006 H8 source is not
authenticated as a complete common tuple. No finite-time unforced
Navier--Stokes singularity is constructed or proved.
